\chapter{Radiation}
\thispagestyle{fancy}

\section{Introduction}

Radiation is the emission and propagation of energy through space or matter. It is a fundamental phenomenon of nature, occurring across an enormous range of scales, from the transmission of radio signals and visible light to the energetic particles produced by nuclear reactions and cosmic events. Radiation serves as one of the primary means by which energy is transferred throughout the universe.

Radiation may take the form of electromagnetic waves, such as radio waves, microwaves, infrared radiation, visible light, ultraviolet radiation, X-rays, and gamma rays, or it may consist of particles, including electrons, protons, neutrons, and atomic nuclei. Regardless of its form, radiation is characterized by its ability to transport energy and momentum away from its source.

The interaction of radiation with matter governs countless natural and technological processes. These interactions determine how stars shine, how living organisms perceive light, how materials are analyzed, and how energy can be generated, transmitted, and measured. Depending on its energy and properties, radiation may pass through matter, be absorbed, scattered, or induce physical and chemical changes within the material it encounters.

The study of radiation lies at the intersection of physics, chemistry, biology, astronomy, and engineering. Understanding its origins, behavior, and effects is essential for fields ranging from medical imaging and radiation therapy to telecommunications, nuclear science, environmental monitoring, and space exploration.

This chapter examines the topic of radiation, which will end up being as broad or as focused as the information I happen to be studying at the time and where that leads me.








\section{Discovery}

Modern history records that Henri Bacquerel first discovered \keyword{radioactivity} \cite{Strom Advanced General Chemistry} at the turn of the 20th century.

\begin{defn}[]{Radioactivity}
	\keyword{Radioactivity} is ``the property possessed by some elements (such as uranium) or isotopes (such as carbon 14) of spontaneously emitting energetic particles (such as electrons or alpha particles) by the disintegration of their atomic nuclei \cite{Dictionary}.''
\end{defn}

Early work with radioactivity involved experiments that showed some radiation could pass through aluminum foil - while others couldn't. It was determined by these that there were three basic types of radiation, which were named alpha, beta, and gamma radiation. A more modern and current understanding is that the \keyword{alpha particle} ($\alpha$) is a helium nucleus (two protons and two neutrons), a \keyword{beta particle} ($\beta$) is a high speed electron, and \keyword{gamma rays} ($\gamma$) are high energy forms of electromagnetic waves (light).